\scanpage{13}
l'injectif indécomposable

Ces conditions impliquent la suivante

c) Pour tout hom $u\colon R\mathbin{-}R'$ surjectif de modules de type fini sur un $\mathcal O_{\deleted{\unsure{M|A}}X,x}$ $(x\in U)$, tel que $\unsure{\operatorname{Hom}|\operatorname{Rhom}}(R',\Omega_x/(f_1\ldots f_n)\Omega_x)\to\operatorname{Hom}(R,\Omega_x/(f_1\ldots f_n)\Omega_x)$ soit un isomorphisme, \inserted{pour tt syst. de paramètres $f_1\ldots f_n$ de $\mathcal O_{X,x}$,} \deleted{$\unsure{e|c}$}$u$ est \unsure{lui $\overline{\mathrm m}$|lui-m} un isomorphisme

% Several diagonals cancel the Corollaire 11 draft below. Horizontal strokes and local rewrites are distinguished inside the cancellation.
\deleted{\uline{Corollaire 11} \deleted{Supposons \unsure{que}} \inserted{Supposons que} $U,\Omega$ soient comme dans Cor. 10. Soit $x\in U$, et soit $u\colon R\mathbin{-}R'$ un hom \inserted{\uline{surjectif}} de modules de type fini sur $\mathcal O_{X,x}$, \deleted{qui soit} tel que $R$ et $R'$ satisfassent la condition suivante pour les foncteurs $\Phi_R(M)=\deleted{\operatorname{Hom}(R,\unsure{D\otimes M|D M|\Omega\otimes M})}\inserted{\operatorname{Hom}(R,M)}$ et $\Phi_{R'}(M)=\deleted{\operatorname{Hom}(R',\unsure{D\otimes M|D M|\Omega\otimes M})}\operatorname{Hom}(R',M)$ :}

% slanted
\margin{à expliciter d'une autre façon \tikz{\draw (0,0)--(0,1.1);}}

\deleted{a)}
\[
\deleted{\deleted{\unsure{\Phi_R|\Phi}}\Phi(\Omega)\otimes A/(\deleted{\unsure{f}}f_1,\ldots f_n)\longrightarrow\Phi(\Omega\otimes A/(f_1\ldots f_n))}
\]
\deleted{est un isomorphisme pour \deleted{\unsure{tt}} \deleted{système de paramètres} \inserted{\unsure{suite|suites} $A$-régulière} $(f_1\ldots f_n)$ de $\mathcal O_{X,x}$ (où $\Phi$ est $\Phi_R$ ou $\Phi_{R'}$)}

\deleted{b)}
\[
\deleted{\deleted{\unsure{\operatorname{Hom}(R|\operatorname{Hom}}}\Phi_{R'}(\Omega)\xrightarrow{\sim}\Phi_R(\Omega)}
\]

\deleted{Sous ces conditions, $u$ est un \uline{isom.}}

% archive sheet number: 12
